\begin{abstract}
  While model merging and low-rank adaptation (LoRA) enable resource-efficient multi-task specialization, their behavior under real-world post-training quantization (PTQ) remains poorly understood. In this paper, we deconstruct the ``Re-Quantization Silence'': the phenomenon where post-hoc low-bit quantization (e.g., 4-bit) of merged models truncates task-specific adapter updates to zero. To systematically audit this behavior, we present a rigorous Multi-Axial Re-Quantization Auditing (RQA) framework. We expose a crucial bifurcation: while naive re-quantization is nearly lossless under standard per-channel grids (dropping only 0.15\% to 0.30\% accuracy), it collapses catastrophically under aggressive per-tensor constraints, losing up to 8.6\% mean accuracy. To mitigate these drops, we evaluate two proposed interventions: Scale-Adaptive Weight Shifting (SAWS), a data-free scaling method, and Quantization-Aware Adapter Coefficient Search (QA-ACS), a test-time optimization method. Crucially, we conduct a self-critical deconstruction of these methods: we expose a fundamental ``representation scale preservation dilemma'' in SAWS, showing that true mathematical scale preservation collapses base representations and that its success is driven by selective task-vector boosting. We also examine the robustness and limits of test-time optimization (QA-ACS and AdaMerging), demonstrating that they achieve the best overall accuracy under standard per-channel grids, but can exhibit minor local instability under aggressive per-tensor noise, which is completely resolved via basic coefficient regularization. By laying bare these empirical and mathematical realities, this work advocates for a more transparent, self-critical, and deployment-aware approach to model-merging evaluation.
\end{abstract}